\section{Language Overview}

\plat{} programs use the \code{.plat} extension and are executed by the
\code{platlang} interpreter. The current repository version is \code{0.1.0}.
The language is dynamically typed and currently includes numbers, strings,
booleans, \code{niks}, and mutable \code{portefeuil} tables. Runtime type names in
the design notes include \code{nommer}, \code{teks}, \code{woar}, \code{niks},
and \code{portefeuil}. Functions are global declarations rather than first-class
runtime values.

The basic syntax is statement-oriented. Variables are declared with
\code{loat}. Functions are declared with \code{funksie} and return with
\code{trok}. Conditional statements use \code{es}, optional \code{angesj}, and
\code{enj}. The language includes \code{zolang} loops and half-open numeric
\code{veur} loops. Loop control uses \code{aafbraeke} and \code{euversjlaon}.

\begin{lstlisting}[caption={A recursive Fibonacci example from the repository.},label={lst:fib}]
# Recursive function example.

funksie fib(n):
    es n < 2:
        trok n
    angesj:
        trok fib(n - 1) + fib(n - 2)
    enj
enj

aafdrokke(fib(10))
\end{lstlisting}

Tables are the only composite data structure in the current language design.
They can be used as arrays, maps, records, objects, or namespaces. Array-style
use is zero-based, field access supports both dot and index notation, and keys
must evaluate to strings or numbers.

\begin{lstlisting}[caption={A table example based on the Limburgian-commented examples.},label={lst:tables}]
loat gebroeker = {
    naam = "Alice",
    leefteid = 30,
}

loat nommers = { 10, 20, 30, }
loat veld = "naam"

aafdrokke(gebroeker.naam)
aafdrokke(gebroeker[veld])
aafdrokke(nommers[0])
\end{lstlisting}

The language deliberately omits many features common in larger systems:
classes, inheritance, static typing, generics, macros, exceptions, operator
overloading, first-class functions, closures, nested functions, and module
systems. These omissions should not be read as claims about what programming
languages ought to exclude in general. They are local choices that keep the
case study small enough for the cultural design layer to remain legible.

% TODO: Verify the feature list against a tagged release before publication.
% TODO: Add a compact grammar table or reference appendix once the syntax is
% stable.
